\scanpage{78}
\uline{Conclusion} Les catégories $\underline{\mathrm{Equ}}(C)$, $C\in\mathrm{ob}_0\underline{\underline B}_0$, sont \deleted{des gr-catégories munies de str} \inserted{munies canoniquement d'une structure de} catégorie de Picard \uline{stricte\deleted{s}} \deleted{\uline{\unsure{commutatives}}}. \deleted{Ils sont} Ces cat. de Picard sont \uline{canoniquement} équivalentes entre elles, et peuvent donc s'identifier à une m\textsuperscript{e} catégorie de Picard $B$ (catégorie de Picard de \uline{Brandt}). Donc $\underline B$ opère sur chaque $C$, et pour $C^{\unsure{*|{}}},C'\in\mathrm{ob}_0\underline{\underline B}_0$, la catégorie \deleted{Hom} $\underline{\mathrm{Equ}}(C,C')$ est canoniquement munie d'une structure de tors-catégorie sous $\underline B$. De façon précise,

\[
\begin{aligned}
C'&\longmapsto\underline{\mathrm{Equ}}(C,C')\\
\underline{\underline B}_0&\longrightarrow \text{2-catégorie des tors-catégories}\\
&\hspace{5em}\text{sous }\underline B
\end{aligned}
\]

est une \unsure{équivalence|équivalente} de $2$-catégories.

Si on a un chgt de base

\diagram{\begin{tikzcd}[ampersand replacement=\&,row sep=small,column sep=small]
A\arrow[r]\arrow[d,hook,dash] \& A'\arrow[d,hook,dash]\\
K\arrow[r,hook,dash] \& K'
\end{tikzcd}}

, et si les $C_{A'}$ satisfont les m\textsuperscript{es} conditions que les $C$, alors on a un hom de \deleted{gr} catégories de Picard

\[
\underline B\longrightarrow B'
\]

et l'hom \deleted{\unsure{de} \illegible{2}}

\[
C\longrightarrow C_{A'}
\]

est compatible avec les opérations de $B,B'$\ldots
% archive sheet number: 77

